\section{Related work}
\label{sec:related}

This work sits at the intersection of four literatures.

\textbf{Reservoir computing and QRC.} Classical reservoir computing emphasizes fixed nonlinear dynamics with a trained readout \cite{maass2002real,lukosevicius2009reservoir,tanaka2019recent}. QRC inherits that structure while replacing the reservoir with a quantum dynamical system \cite{fujii2017harnessing,nakajima2019boosting,mujal2021opportunities}. Much of the QRC literature asks what quantum dynamics, noise, dissipation, or many-body effects increase computational power \cite{kubota2022noise,kubota2023temporal,fry2023optimizing,domingo2023taking,sannia2024dissipation,martinez2021dynamical,palacios2024role}. Our question is orthogonal: given a finite shot record from a fixed reservoir, how should that record be organized before the readout is trained?

\textbf{Finite-shot and noise-aware QRC.} The closest finite-shot line is the recent study of finitely sampled QRC by Ahmed, Tennie, and Magri \cite{ahmed2025optimal}, who develop noise-aware training after expectation-value collapse. Our contribution is complementary. We intervene \emph{before} EV collapse by reorganizing the same shot record into grouped views. This is why the empirical section includes both protocol controls (EV, Raw, EV-dup) and stronger finite-shot baselines (EV-NA and Split-NA): the goal is not to claim that grouping replaces estimator-level corrections, but to isolate a separate protocol axis. A full implementation of the Ahmed--Tennie--Magri estimator and its composition with split-ensemble is an important next step; our present EV-NA and Split-NA baselines should therefore be read as a first diagonal proxy rather than as an exhaustive comparator.

\textbf{Measurement error.} The harmful effect of noisy covariates on linear regression is classical \cite{fuller1987measurement,carroll2006measurement}. Raw stacking is exactly the extreme noisy-covariate protocol in our setting, and the reduced forecast amplitude we observe is the expected attenuation phenomenon. The grouped-view surrogate in Section~\ref{sec:methodology} is deliberately written in the language of this literature rather than as a bespoke quantum theorem.

\textbf{High-dimensional ridge and random-matrix conditioning.} When the number of rows is comparable to the feature dimension, ridge prediction depends strongly on conditioning and regularization. This has been analyzed extensively through random-design asymptotics and Marchenko--Pastur theory \cite{marchenko1967distribution,bai2010spectral,hsu2012random,dobriban2018high}. Our theoretical framing follows this standard style: combine one exact protocol transformation with a familiar conditioning lens, then test the implied regime picture empirically. In that sense the paper is theory-backed but intentionally not theorem-heavy.
